\chapter{Literature Review}
\label{chapter:literature_review}

This section reviews existing scientific publications to establish a baseline for data anonymization services within software engineering productivity metrics.

\section{Data Protection Laws}
\label{sec:data_protection_laws}

Finding publications about this topic is very difficult. There are a lot of sources about the \ac{DSGVO} and anonymization and their applications. But most of them are too broad for this thesis or focus on surveillance tools.

\textcite{Christl2021} provides a systematic mapping of digital surveillance and algorithmic management across diverse industries. Through case studies like Amazon and Zalando, he documents a trend toward "digital Taylorism," characterized by real-time tracking that intensifies work and reduces employee autonomy. Ultimately, the study shows how these technologies reinforce power imbalances between employers and labor through pervasive data collection and automated control.

The standard reference work on employee data protection is \textcite{daeubler2017glaeserne}. Written in a practical style and accessible even to non-specialists, this handbook explains the entire field of data protection law. It focuses on employee data protection and the participation rights of works councils and staff councils. It also addresses the challenges posed by AI that are already apparent today, such as those arising in the hiring process or when working with ChatGPT. Another key focus is on the diverse means available to ensure that data protection is effectively enforced in practice. These range from data backup and the activities of the company data protection officer and the supervisory authority to the imposition of substantial fines and numerous claims for damages. The now very extensive case law of the European Court of Justice (ECJ) — which has set many new precedents — is analyzed in relation to these and numerous other questions. As a rule, this case law strengthens the protection of individuals, which is also of great significance for the scope of action available to works councils and staff councils. This shift has not yet been fully recognized everywhere.

In the context of data protection and research, the operationalization of de-identification procedures under the \ac{DSGVO} is addressed by two central German guidelines that provide both technical and legal frameworks. \textcite{schwartmann2022praxisleitfaden} establishes a structured procedural model that distinguishes between theoretical "absolute anonymity" and the more practically relevant "factual anonymization," where re-identification is considered impossible due to the disproportionate effort required in terms of time, cost, and manpower. Complementing this, \textcite{gmdsbvd2024praxishilfe} focuses specifically on the sensitive domain of health data under Art. 9 \ac{GDPR}, highlighting the necessity of aligning national definitions with the broader European legal framework provided by Directive (EU) 2019/1024. While \textcite{schwartmann2022praxisleitfaden} categorizes technical methods—such as randomization (e.g., differential privacy) and generalization (e.g., k-anonymity)—into specific application classes like AI algorithm training or software testing, \textcite{gmdsbvd2024praxishilfe} provides detailed risk assessments for specific attack scenarios, including "singling out," inference, and linkage attacks. Both sources emphasize that anonymity is not a static state but a dynamic process that must be continuously evaluated against the evolving "state of science and technology". Furthermore, they converge on the legal requirement that both anonymization and pseudonymization constitute "processing" activities under the  \ac{GDPR}, therefore requiring a valid legal basis and remaining subject to the controller's accountability obligations.

\section{Server Design}
\label{sec:server_design}

TODO

\section{Identity Management}
\label{sec:identity_management}

The MECOIS project uses the tool SortinhHat from the toolset GrimoireLab by CHAOSS for managing identities.

\textcite{duenas2021grimoirelab} present GrimoireLab, an integrated and modular open-source toolset developed to facilitate the comprehensive retrieval, storage, and visualization of data from software repositories. This framework addresses a critical gap in software analytics by providing automated support for approximately 30 different data sources, including version control systems like Git, issue trackers such as Jira and Bugzilla, and communication platforms like Slack and mailing lists. Central to its architecture is a multi-stage pipeline that employs specialized components: Perceval for uniform data retrieval via JSON documents, Graal for detailed source code analysis, and SortingHat for advanced identity management and cross-platform contributor merging. Furthermore, GrimoireLab utilizes a distinct storage model that maintains "raw" copies of original data alongside "enriched" indexes optimized for visualization, a strategy that significantly enhances data maintainability, traceability, and research reproducibility. By providing a mature, field-tested framework for both academic research and industrial application, the authors demonstrate that GrimoireLab reduces the manual effort required for empirical software engineering studies while enabling the longitudinal tracking of developer activity and project health.

SortingHat is the specialized identity management component within the GrimoireLab framework, specifically designed to address the challenge of contributor fragmentation across diverse software repositories. Because developers frequently utilize different usernames, aliases, and email addresses across platforms such as Git, Slack, and Jira, SortingHat provides the essential infrastructure to unify these disparate identities into unique, individual profiles. The component operates by maintaining a dedicated relational database that stores identity data, affiliation tags, and the original source of each identity to ensure research traceability. To maintain high data quality, SortingHat utilizes a conservative algorithmic approach to merging—avoiding false positives caused by shared generic accounts—which is further complemented by HatStall, a web-based interface that allows researchers to manually curate identities and manage organizational affiliations. This process of identity unification is fundamental to the broader GrimoireLab pipeline; it allows the analytics engine to enrich data points with a unique contributor \ac{UUID}, thereby enabling the longitudinal tracking of developer trajectories and cross-platform activity metrics that would otherwise remain siloed. \autocite{duenas2021grimoirelab}
